% !TEX root = complete_curved_transport.tex
\documentclass[a4paper,11pt]{ltjsarticle}
\usepackage[margin=23mm,headheight=16pt]{geometry}
\usepackage{amsmath,amssymb,amsthm,mathtools}
\usepackage[haranoaji]{luatexja-preset}
\usepackage[unicode,hidelinks]{hyperref}
\usepackage{fancyhdr}
\newtheorem{theorem}{定理}[section]
\newtheorem{lemma}[theorem]{補題}
\newtheorem{corollary}[theorem]{系}
\theoremstyle{remark}\newtheorem{remark}[theorem]{注意}
\renewcommand{\proofname}{証明}
\numberwithin{equation}{section}
\setlength{\emergencystretch}{3em}
\allowdisplaybreaks
\pagestyle{fancy}\fancyhf{}
\fancyhead[L]{\small 完備曲率付き移送と形式的MF比較}
\fancyhead[R]{\small 2026-11-12 / B-v43}
\fancyfoot[C]{\thepage}
\makeatletter
\def\input@path{{papers/support/}}
\AtBeginDocument{\let\originalref\ref
\renewcommand{\ref}[1]{\ifcsname r@#1\endcsname\originalref{#1}\else\texttt{\detokenize{#1}}\fi}}
\makeatother
\title{完備曲率付き移送と形式的MF比較\\\large 固定三対象の比較条件の構成}
\author{}\date{2026年11月12日 / B-v43・証明抜粋}
\begin{document}
\maketitle
\noindent 本文外への参照は原稿のラベル名で表示する。\\
正本：\texttt{uedakazushi/nc-cubic}、\texttt{papers/paper\_b.tex}。
本抜粋の組版確認は、論文B全文の組版や全証明の形式検証とは異なる。
\input{papers/reconstruction/complete_curved_transport}
\begin{thebibliography}{9}
\bibitem{PKontsevich}
M. Kontsevich, \emph{Deformation quantization of Poisson manifolds},
Letters in Mathematical Physics \textbf{66} (2003), 157--216.
\url{https://arxiv.org/abs/q-alg/9709040}.
著者最終稿\url{https://www.ihes.fr/~maxim/TEXTS/DefQuant_final.pdf}の
Theorem 4.6, Section 4.6.2, Sections 6.2--6.4, Theorem 6.3。
\bibitem{CT}
A. C\u{a}ld\u{a}raru and J. Tu, \emph{Curved $A_\infty$ algebras and Landau--Ginzburg models},
New York Journal of Mathematics \textbf{19} (2013), 305--342.
Section 2, Remark 2.8, Theorem 2.12.
\url{https://nyjm.albany.edu/j/2013/19-16v.pdf}.
\end{thebibliography}
\end{document}
